% Options for packages loaded elsewhere
\PassOptionsToPackage{unicode}{hyperref}
\PassOptionsToPackage{hyphens}{url}
\documentclass[
  ignorenonframetext,
]{beamer}
\newif\ifbibliography
\usepackage{pgfpages}
\setbeamertemplate{caption}[numbered]
\setbeamertemplate{caption label separator}{: }
\setbeamercolor{caption name}{fg=normal text.fg}
\beamertemplatenavigationsymbolsempty
% remove section numbering
\setbeamertemplate{part page}{
  \centering
  \begin{beamercolorbox}[sep=16pt,center]{part title}
    \usebeamerfont{part title}\insertpart\par
  \end{beamercolorbox}
}
\setbeamertemplate{section page}{
  \centering
  \begin{beamercolorbox}[sep=12pt,center]{section title}
    \usebeamerfont{section title}\insertsection\par
  \end{beamercolorbox}
}
\setbeamertemplate{subsection page}{
  \centering
  \begin{beamercolorbox}[sep=8pt,center]{subsection title}
    \usebeamerfont{subsection title}\insertsubsection\par
  \end{beamercolorbox}
}
% Prevent slide breaks in the middle of a paragraph
\widowpenalties 1 10000
\raggedbottom
\AtBeginPart{
  \frame{\partpage}
}
\AtBeginSection{
  \ifbibliography
  \else
    \frame{\sectionpage}
  \fi
}
\AtBeginSubsection{
  \frame{\subsectionpage}
}
\usepackage{iftex}
\ifPDFTeX
  \usepackage[T1]{fontenc}
  \usepackage[utf8]{inputenc}
  \usepackage{textcomp} % provide euro and other symbols
\else % if luatex or xetex
  \usepackage{unicode-math} % this also loads fontspec
  \defaultfontfeatures{Scale=MatchLowercase}
  \defaultfontfeatures[\rmfamily]{Ligatures=TeX,Scale=1}
\fi
\usepackage{lmodern}
\ifPDFTeX\else
  % xetex/luatex font selection
\fi
% Use upquote if available, for straight quotes in verbatim environments
\IfFileExists{upquote.sty}{\usepackage{upquote}}{}
\IfFileExists{microtype.sty}{% use microtype if available
  \usepackage[]{microtype}
  \UseMicrotypeSet[protrusion]{basicmath} % disable protrusion for tt fonts
}{}
\makeatletter
\@ifundefined{KOMAClassName}{% if non-KOMA class
  \IfFileExists{parskip.sty}{%
    \usepackage{parskip}
  }{% else
    \setlength{\parindent}{0pt}
    \setlength{\parskip}{6pt plus 2pt minus 1pt}}
}{% if KOMA class
  \KOMAoptions{parskip=half}}
\makeatother
\usepackage{longtable,booktabs,array}
\newcounter{none} % for unnumbered tables
\usepackage{calc} % for calculating minipage widths
\usepackage{caption}
% Make caption package work with longtable
\makeatletter
\def\fnum@table{\tablename~\thetable}
\makeatother
\setlength{\emergencystretch}{3em} % prevent overfull lines
\providecommand{\tightlist}{%
  \setlength{\itemsep}{0pt}\setlength{\parskip}{0pt}}
\usepackage{bookmark}
\IfFileExists{xurl.sty}{\usepackage{xurl}}{} % add URL line breaks if available
\urlstyle{same}
\hypersetup{
  hidelinks,
  pdfcreator={LaTeX via pandoc}}

\author{\texorpdfstring{}{}}
\date{}

\begin{document}

\begin{frame}{Appendix B: Tooling and Infrastructure
\label{sec:tooling}}
\protect\phantomsection\label{appendix-b-tooling-and-infrastructure}
The practical utility of a computational meta-analysis depends on robust
tooling at each pipeline stage: assertion extraction, modeling and
simulation, knowledge graph infrastructure, and quality assurance.

\begin{block}{LLM-Based Assertion Extraction}
\protect\phantomsection\label{llm-based-assertion-extraction}
Extracting structured assertions from unstructured text is the most
labor-intensive component of knowledge graph construction. Manual
annotation produces high-quality results but does not scale to corpora
of thousands of papers---a constraint demonstrated by Knight et
al.~\citep{knight2022fep}, whose systematic literature analysis of FEP
and Active Inference publications required manual coding of structural,
visual, and mathematical features for hundreds of annotated papers. We
implement a hybrid approach: LLMs perform initial extraction, with human
review for validation and correction.

Our extraction pipeline deploys a locally hosted LLM through Ollama
\citep{ollama2024}. Each paper's abstract is assessed against the eight
hypothesis definitions in a structured prompt requesting a JSON array of
assessments. Unlike keyword matching, which detects only topical terms,
the LLM evaluates the \emph{semantic relationship} between a paper's
claims and each hypothesis. Papers critiquing the FEP correctly receive
``contradicts'' assessments for FEP Universality (H1), while methodology
tutorials receive ``neutral'' assessments reflecting their pedagogical
character. Detailed prompt engineering, schemas, and failure modes are
documented in the supplementary extraction pipeline (see the
\hyperref[sec:extraction_pipeline]{extraction pipeline section}).
\end{block}

\begin{block}{Software Ecosystem}
\protect\phantomsection\label{software-ecosystem}
The Active Inference community has developed a rapidly growing ecosystem
of open-source tools spanning multiple programming languages, inference
paradigms, and application domains. This section provides a
comprehensive survey of publicly available implementations as of early
2026, organized by functional category. We emphasize tools with
accessible source code, as open-source availability is a prerequisite
for reproducibility and community-driven validation.

\begin{block}{General-Purpose Frameworks}
\protect\phantomsection\label{general-purpose-frameworks}
Four general-purpose frameworks dominate the landscape, collectively
covering discrete, continuous, and real-time inference:

\textbf{pymdp.} The pymdp library \citep{heins2022pymdp} provides a
Python implementation of active inference for discrete state-space
POMDPs, supporting message passing on factor graphs, policy inference
via expected free energy, and hierarchical generative models. It has
become the standard entry point for algorithm development and the most
widely forked AIF repository.

\textbf{SPM.} The SPM package (Wellcome Centre for Human Neuroimaging)
includes MATLAB implementations of Dynamic Causal Modeling and
variational Bayesian inference under the FEP. It remains the reference
implementation for neuroimaging applications and houses the original
Friston-group POMDP scripts.

\textbf{RxInfer.jl.} RxInfer is a Julia package for reactive
message-passing-based Bayesian inference, supporting real-time and
streaming inference suitable for robotics and online learning. Version
4.0.0 (early 2025) \citep{rxinfer2025} introduced projected constraints
and adaptive inference optimized for dynamic data streams and autonomous
systems. The RxInfer ecosystem includes extensive tutorials covering
Bayesian linear regression, hidden Markov models, Kalman filtering,
Gaussian process regression, hierarchical Gaussian filters, nonlinear
sensor fusion, and active inference mountain car control, available at
the \href{https://reactivebayes.github.io/RxInfer.jl/stable/}{official
documentation} and the \href{https://learnableloop.com/}{Learnable Loop}
tutorial portal.

\textbf{Cpp-AIF.} The Cpp-AIF header-only C++ library
\citep{gregoretti2023cppaif} implements active inference for discrete
POMDPs with multicore parallelization of the most demanding
computational kernels---multidimensional inner products for expected
free energy computation and state estimation. By abstracting the
mathematical details behind a high-level API, Cpp-AIF targets embedded
systems and performance-critical applications where Python overhead is
prohibitive.

\textbf{FEPS.} Free Energy Projective Simulation \citep{pazem2024feps}
combines active inference with interpretable graphical policy
representations, enabling agents to plan via expected free energy while
exposing decision logic as human-readable policy graphs. FEPS targets
interpretable reinforcement learning tasks where black-box deep agents
are undesirable---behavioral biology, clinical decision support, and
safety-critical robotics.
\end{block}

\begin{block}{Deep Active Inference}
\protect\phantomsection\label{deep-active-inference}
Scaling active inference beyond tabular POMDPs to high-dimensional
observation spaces requires neural network function approximators. A
growing body of deep active inference implementations explores this
direction:

The foundational deep AIF agent of Fountas et
al.~\citep{fountas2020deep} introduced Monte-Carlo tree search over
learned latent spaces, achieving non-trivial Atari performance.
Millidge's DeepActiveInference extended this to continuous control with
backpropagation-based world models \citep{millidge2020deep}. Champion's
Branching-Time Active Inference (BTAI\_3MF) and its deep variant
(Deep\_BTAI\_3MF) implement tree-structured planning under the free
energy objective, scaling active inference to partially observable
environments with multi-step lookahead \citep{champion2021realizing}.
Most recently, AXIOM \citep{heins2025axiom} achieves competitive Atari
game performance using expanding object-centric world models, learning
in minutes rather than hours---a landmark result for scalability.
\end{block}

\begin{block}{Predictive Coding and Neural Generative Coding}
\protect\phantomsection\label{predictive-coding-and-neural-generative-coding}
Predictive coding provides the core computational mechanism linking
active inference to neuroscience. Several implementations offer
accessible entry points:

\textbf{ngc-learn.} The Neural Generative Coding library (ngc-learn
v3.0, JAX-based) provides a framework for simulating
neurobiologically-plausible systems using predictive coding circuits,
Hebbian learning, and spike-based dynamics. It supports constructing
arbitrary neural generative models without backpropagation, directly
instantiating the FEP's prediction-error minimization at the circuit
level.

\textbf{Active Neural Generative Coding (ANGC).} ANGC implements a form
of active inference using paired predictive coding circuits---an
actor/policy circuit and a world/transition model---that co-evolve
across episodes without backpropagation. The agent decomposes behavior
into epistemic foraging (uncertainty reduction) and instrumental
(reward-seeking) terms, operating with sparse rewards where classical
DQN requires dense reward engineering.

\textbf{Predictive Coding \(\approx\) Backprop.} Millidge et
al.~demonstrate that predictive coding networks can approximate
backpropagation along arbitrary computational graphs
\citep{millidge2022predictive}, providing a biologically plausible
alternative to gradient descent. The
\href{https://github.com/BerenMillidge/PredictiveCodingBackprop}{PredictiveCodingBackprop}
repository provides the reference implementation.
\end{block}

\begin{block}{Benchmarking Progress}
\protect\phantomsection\label{benchmarking-progress}
The scalability gap between AIF and deep reinforcement learning has been
a central limitation of the tools domain. Recent work demonstrates
significant progress on two fronts. First, AXIOM \citep{heins2025axiom}
outperforms state-of-the-art model-based deep RL agents including
DreamerV3 on the Gameworld 10k benchmark, while using substantially
smaller model sizes; its object-centric scene decomposition enables
sample-efficient learning from structured representations rather than
raw pixel memorization. Second, variational message passing formulations
\citep{champion2021realizing} connect EFE decomposition---into risk,
ambiguity, epistemic (information-seeking), and instrumental
(goal-reaching) components---to practical planning algorithms, advancing
the theoretical justification for EFE-based policy selection (H2).
Separately, Friston et al.~\citep{friston2025active} introduce structure
learning via Bayesian Model Reduction as a principled approach to
artificial reasoning under active inference.
\end{block}

\begin{block}{Comprehensive Open-Source Tool Survey}
\protect\phantomsection\label{comprehensive-open-source-tool-survey}
The following table catalogs the principal open-source Active Inference
implementations surveyed, organized by functional category. For each
tool we list the primary language, the application domain, and the
associated publication or repository. This table is intended as a
navigational resource for researchers seeking existing implementations
relevant to specific hypotheses (H1--H8) or application domains
(A1--C5).

\begin{center}
\small
\begin{longtable}{p{3.2cm} p{1.6cm} p{4.4cm} p{2.5cm}}
\hline
\textbf{Tool / Repository} & \textbf{Lang.} & \textbf{Description} & \textbf{Paper / Source} \\
\hline
\endfirsthead
\hline
\textbf{Tool / Repository} & \textbf{Lang.} & \textbf{Description} & \textbf{Paper / Source} \\
\hline
\endhead
\multicolumn{4}{c}{\textit{General-Purpose Frameworks}} \\
\hline
pymdp & Python & Discrete POMDP active inference; factor graphs, hierarchical models & \cite{heins2022pymdp} \\
SPM & MATLAB & DCM, variational Bayes; neuroimaging reference implementation & \cite{friston2017active} \\
RxInfer.jl & Julia & Reactive message passing; real-time streaming Bayesian inference & \cite{rxinfer2025} \\
Cpp-AIF & C++ & Header-only POMDP AIF library with multicore parallelization & \cite{gregoretti2023cppaif} \\
FEPS & Python & EFE on interpretable policy graphs; projective simulation & \cite{pazem2024feps} \\
ActivPynference & Python & Discrete AIF with factor-graph message passing; educational focus & — \\
pypc & Python & Predictive coding inference engine for continuous models & — \\
ActiveInferAnts & Rust & Rust-native AIF framework with WASM compilation target & — \\
\hline
\multicolumn{4}{c}{\textit{Deep Active Inference}} \\
\hline
deep-active-inference-mc & Python & Monte-Carlo tree search in learned latent spaces; Atari & \cite{fountas2020deep} \\
DeepActiveInference & Python & Continuous deep AIF with backprop-based world models & \cite{millidge2020deep} \\
BTAI\_3MF & Python & Branching-time AIF with multi-step tree planning & \cite{champion2021realizing} \\
Deep\_BTAI\_3MF & Python & Deep neural variant of BTAI with learned state spaces & \cite{champion2021realizing} \\
OO-BTAI\_3MF & Python & Object-oriented BTAI variant for structured environments & — \\
AXIOM & Python & Object-centric world models; Atari in minutes, beats DreamerV3 & \cite{heins2025axiom} \\
Deep-AIF-POMDPs & Python & Deep AIF for partially observable MDPs & — \\
Homing-Pigeon & Python & Navigation agent using deep active inference & — \\
active-inference (Voostrum) & Python & Continuous deep AIF with learned generative models & arXiv:2406.07726 \\
\hline
\multicolumn{4}{c}{\textit{Predictive Coding \& Neural Generative Coding}} \\
\hline
ngc-learn & Python/JAX & Neurobiological simulation; predictive coding circuits, Hebbian learning & — \\
ANGC & Python & Backprop-free AIF agent with paired PC circuits & AAAI 2022 \\
PredictiveCodingBackprop & Python & Predictive coding approximates backprop on arbitrary graphs & \cite{millidge2022predictive} \\
Supervised-Predictive-Coding & Python & Supervised learning via hierarchical predictive coding & — \\
predcoding & Python & Minimal predictive coding implementation & — \\
pybrid & Python & Hybrid predictive coding and active inference library & — \\
nmpassing & Python & Neural message passing for PC networks & — \\
\hline
\multicolumn{4}{c}{\textit{Neuroscience, Embodied \& Biological}} \\
\hline
allostasis & Python & Allostatic regulation via AIF; interoceptive inference & bioRxiv:2021.02.16 \\
ants & Python & Ant foraging simulation with stigmergic AIF agents & \cite{heins2024collective} \\
Reward\_Bases & Python & Reward-basis function representations under AIF & bioRxiv:2022.04.14 \\
action-oriented & Python & Action-oriented predictive processing models & \cite{tschantz2020action} \\
Biofirm & Python & Bioregional stewardship via organizational AIF & — \\
bayesian-mechanics-sdes & Python & Bayesian mechanics: SDE simulations of Markov blanket dynamics & arXiv:2206.02629 \\
reverse\_engineering & MATLAB & Reverse-engineering neural dynamics under the FEP & — \\
\hline
\multicolumn{4}{c}{\textit{Multi-Agent \& Social Dynamics}} \\
\hline
opinion\_dynamics & Python & Opinion dynamics and belief formation via AIF & — \\
network-actinf & Python & Network-level active inference with coupled agents & — \\
Variational-Capsule-Routing & Python & Capsule networks with variational inference routing & AAAI 2020 \\
Active-Inference-Successor & Python & Successor representations under active inference & — \\
\hline
\multicolumn{4}{c}{\textit{Domain-Specific Applications}} \\
\hline
adaptive\_aif\_agents\_fl & Python & Adaptive AIF agents for federated learning & arXiv:2410.09099 \\
smartville & Python & IoT smart building control via AIF under partial observability & TechRxiv 2025 \\
FEP\_Blorpomon & Python & Game-theoretic AIF agent demonstration & — \\
MountainCarAI & Python & Mountain car control via active inference & — \\
rl-inference & Python & Bridging RL and active inference policy selection & arXiv:2002.12636 \\
EFE-GLean & Python & Expected free energy with generalized learning & Entropy 2025 \\
EFEasVFE & Julia & EFE reformulated as variational free energy & — \\
Robust-FE-Minimization & Python & Robust decision-making via free energy minimization & arXiv:2503.13223 \\
\hline
\multicolumn{4}{c}{\textit{Tutorials \& Educational Resources}} \\
\hline
Active-Inference-from-Scratch & Python & Step-by-step AIF implementation tutorial & — \\
IC2S2-AIF-Tutorial & Python & Computational social science AIF tutorial & — \\
julia4ta tutorials (9x10--12) & Julia & RxInfer-based AIF agent tutorials & — \\
ActInf Textbook Colab & Python & Interactive notebooks for \cite{parr2022active} & — \\
deep\_aif\_workshop & Python & Workshop materials for deep active inference & — \\
AdaptiveResonance.jl & Julia & Adaptive resonance theory models in Julia & — \\
\hline
\end{longtable}
\end{center}
\end{block}

\begin{block}{Comparative Feature Matrix}
\protect\phantomsection\label{comparative-feature-matrix}
{\def\LTcaptype{none} % do not increment counter
\begin{longtable}[]{@{}
  >{\raggedright\arraybackslash}p{(\linewidth - 12\tabcolsep) * \real{0.1429}}
  >{\raggedright\arraybackslash}p{(\linewidth - 12\tabcolsep) * \real{0.1429}}
  >{\raggedright\arraybackslash}p{(\linewidth - 12\tabcolsep) * \real{0.1429}}
  >{\raggedright\arraybackslash}p{(\linewidth - 12\tabcolsep) * \real{0.1429}}
  >{\raggedright\arraybackslash}p{(\linewidth - 12\tabcolsep) * \real{0.1429}}
  >{\raggedright\arraybackslash}p{(\linewidth - 12\tabcolsep) * \real{0.1429}}
  >{\raggedright\arraybackslash}p{(\linewidth - 12\tabcolsep) * \real{0.1429}}@{}}
\toprule\noalign{}
\begin{minipage}[b]{\linewidth}\raggedright
Feature
\end{minipage} & \begin{minipage}[b]{\linewidth}\raggedright
pymdp
\end{minipage} & \begin{minipage}[b]{\linewidth}\raggedright
SPM
\end{minipage} & \begin{minipage}[b]{\linewidth}\raggedright
RxInfer.jl
\end{minipage} & \begin{minipage}[b]{\linewidth}\raggedright
Cpp-AIF
\end{minipage} & \begin{minipage}[b]{\linewidth}\raggedright
FEPS
\end{minipage} & \begin{minipage}[b]{\linewidth}\raggedright
ngc-learn
\end{minipage} \\
\midrule\noalign{}
\endhead
\textbf{Language} & Python & MATLAB & Julia & C++ & Python &
Python/JAX \\
\textbf{State Spaces} & Discrete & Discrete + Continuous & Continuous
(factor graphs) & Discrete & Discrete & Continuous \\
\textbf{Inference} & Message passing & Variational Bayes & Reactive
message passing & EFE + state estimation & EFE on policy graphs &
Predictive coding \\
\textbf{Deep AIF} & Partial & No & Via custom factors & No & No
(interpretable) & Yes (neural circuits) \\
\textbf{Real-time} & No & No & Yes (streaming) & Yes (multicore) & No &
No \\
\textbf{Hierarchical} & Yes & Yes (DCM) & Yes & Yes & No & Yes \\
\textbf{GPU} & No & No & No & CPU (multicore) & No & Yes (JAX) \\
\textbf{License} & MIT & GPL & MIT & MIT & MIT & BSD-3 \\
\textbf{Primary Use} & Research prototyping & Neuroimaging & Robotics /
online learning & Embedded systems & Interpretable RL & NeuroAI
simulation \\
\bottomrule\noalign{}
\end{longtable}
}

The complementary strengths across these packages reflect a fragmented
but maturing ecosystem. The survey reveals several notable patterns: (1)
Python dominates (\textasciitilde75\% of implementations), with Julia
emerging as the preferred alternative for performance-critical
applications; (2) discrete POMDP implementations outnumber continuous
variants by approximately 3:1, reflecting pymdp's community influence;
(3) deep active inference implementations are concentrated in a small
number of research groups (Champion, Millidge, Fountas, Heins),
suggesting high barriers to entry; (4) multi-agent and social AIF
implementations remain sparse relative to single-agent tools; and (5)
domain-specific applications (IoT, federated learning, smart buildings)
represent the newest and fastest-growing category, aligning with the
temporal growth patterns observed in the C-domain (applied) subfields.
The variational free energy foundations shared by Active Inference and
Energy-Based Models (EBMs)---including Helmholtz machines
\citep{dayan1995helmholtz}, Boltzmann machines
\citep{hinton2002training}, and variational autoencoders
\citep{kingma2014auto}---suggest that interoperability with mainstream
deep generative modeling frameworks (PyTorch, JAX) could bridge these
parallel research programs.
\end{block}
\end{block}

\begin{block}{Knowledge Graph Infrastructure}
\protect\phantomsection\label{knowledge-graph-infrastructure}
Our knowledge graph uses an RDF-compatible schema deployable on standard
semantic web infrastructure. The nanopublication model
\citep{groth2010anatomy, kuhn2016decentralized} provides a principled
atomic unit of scientific evidence: each nanopublication packages a
single assertion (e.g., ``Paper X supports Hypothesis Y'') with explicit
provenance and publication metadata in four named RDF graphs (Head,
Assertion, Provenance, Publication Info). This structure satisfies the
FAIR data principles by design: nanopublications are \textbf{F}indable
via URI-based identification, \textbf{A}ccessible through standard RDF
protocols, \textbf{I}nteroperable via W3C-standard TriG serialization,
and \textbf{R}eusable with explicit provenance and CC0 licensing. The
full RDF schema and a TriG serialization example are presented in the
\hyperref[sec:methods_kg]{methodology} and
\hyperref[sec:appendix_rdf]{Appendix~A.5}.

The engineering trade-offs among the three deployment options are
straightforward:

\textbf{Nanopublication servers} provide decentralized,
content-addressed storage. The pipeline writes nanopublications in two
forms: JSON Lines (for incremental checkpointing and tooling) and
RDF/TriG per the \href{https://nanopub.net/}{nanopublication standard}
(Assertion, Provenance, Publication Info), suitable for the
nanopublication network and FAIR deployment. The recent release of
nanopub-js v0.1.0 \citep{kuhn2026nanopubjs}---a JavaScript library
enabling browser-based creation, signing, and querying of
nanopublications---opens the possibility of community-contributed
assertions directly from web interfaces, lowering the barrier to
participatory evidence curation. Future integration with Trusty URIs
\citep{kuhn2014trusty} would provide cryptographic content verification
and persistent identifiers for each nanopublication.

\textbf{RDF stores} (e.g., Apache Jena Fuseki, Blazegraph, Oxigraph)
enable SPARQL queries such as ``find all papers supporting hypothesis H
published after 2020 in the neuroscience domain (C1).'' The cost is
operational overhead and query latency.

\textbf{Property graph databases} (e.g., Neo4j) prioritize traversal
performance for path queries and community detection, at the expense of
semantic web compatibility.

The \href{http://activeinference.org/ontology/}{Active Inference
Ontology namespace} ensures integration with external ontologies and
linked data resources.
\end{block}

\begin{block}{Multi-Level Quality Assurance}
\protect\phantomsection\label{multi-level-quality-assurance}
Quality assurance operates at four levels.

\begin{block}{Assertion-Level Validation}
\protect\phantomsection\label{assertion-level-validation}
Assertions below a configurable confidence threshold (default 0.5) are
flagged for review. Inter-annotator agreement (\(\kappa\)) is computed
when multiple annotators assess the same paper.
\end{block}

\begin{block}{Graph-Level Consistency Checks}
\protect\phantomsection\label{graph-level-consistency-checks}
Consistency checks verify that all nodes link to valid targets and no
orphan nodes exist. Coverage metrics track the proportion of annotated
papers.
\end{block}

\begin{block}{Score-Level Unit Testing}
\protect\phantomsection\label{score-level-unit-testing}
Hypothesis scoring is validated through unit tests with synthetic data
verifying boundary conditions (all-support → +1, all-contradict → −1,
balanced → 0). Sensitivity analysis varies confidence thresholds and
citation weighting.
\end{block}

\begin{block}{Pipeline-Level Test Coverage}
\protect\phantomsection\label{pipeline-level-test-coverage}
Test-driven development enforces 90\% minimum code coverage on project
modules and 60\% on shared infrastructure, with real data and
computation (no mocking).
\end{block}

\begin{block}{Quality Thresholds}
\protect\phantomsection\label{quality-thresholds}
{\def\LTcaptype{none} % do not increment counter
\begin{longtable}[]{@{}llll@{}}
\toprule\noalign{}
Level & Metric & Threshold & On Failure \\
\midrule\noalign{}
\endhead
Assertion & Confidence & \(\geq 0.5\) & Flag for review \\
Assertion & Inter-annotator \(\kappa\) & \(\geq 0.6\) & Re-annotate \\
Graph & Orphan node ratio & \(= 0\) & Reject build \\
Graph & Corpus coverage & \(\geq 80\%\) & Warning \\
Score & Boundary tests & All pass & Block release \\
Pipeline & Code coverage & \(\geq 90\%\) & Block merge \\
Pipeline & Test pass rate & \(100\%\) & Block release \\
\bottomrule\noalign{}
\end{longtable}
}

The hypothesis evidence results, temporal dynamics of evidence
accumulation, and assertion analysis are presented in the dedicated
hypothesis results section (see the
\hyperref[sec:hypothesis_results]{hypothesis results section}).
\end{block}
\end{block}
\end{frame}

\end{document}
